\cleardoublepage
\chapter{HASIL SURVEI LAPANGAN}
\section{Foto Survei Lokasi 1}
Teks survei lokasi 1.
\section{Foto Survei Lokasi 2}
Teks survei lokasi 2.
\section{Validasi Prototipe dan Diskusi Pengembangan Lanjutan}

\begin{table}[H]
\centering
\begin{tabular}{|p{3cm}|p{11cm}|}
\hline
\textbf{Hari, Tanggal} & Kamis, 21 Juni 2026 \\ \hline
\textbf{Pukul} & 13.30 \\ \hline
\textbf{Narasumber} &
Dr. Roos Meri Silaen.
(Dokter Umum) \\ \hline
\textbf{Tempat} &
KK Farmakologi, Lt. 3 Labtek VII Gedung Sekolah Farmasi ITB \\ \hline
\textbf{Notulensi} &
\begin{enumerate}
    \item Prototipe sistem diinterpretasikan kepada narasumber.
    
    \item Fitur ``Filter Alergi Pribadi'' divalidasi sebagai fitur yang sangat relevan dan krusial, mengingat sifat alergi yang sangat individual.
    
    \item Wawasan mengenai minyak yang bukan termasuk alergen utama. Namun, karakternya dapat berubah ketika bertemu dengan protein (seperti pertemuan dengan ayam menjadi minyak ayam). Hal ini memungkinkan potensi risiko pada individu tertentu yang dianggap sensitif.
    
    \item Disarankan fitur baru untuk risiko laptem dari minyak/lemak dikomersialisasikan sebagai pengingat yang sangat bernilai untuk edukasi dan keamanan pengguna.
    
    \item Sepakat bahwa batasan-batasan sistem (tidak mencakup kontaminasi silang dan transformasi kimia) telah diidentifikasi dengan tepat dalam penelitian.
    
    \item Arah pengembangan selanjutnya yang paling berdampak yaitu adisi modul meningkatkan \textit{avoidability} (menjelaskan penyebab \textit{flag}) dan membangun mekanisme \textit{feedback loop} untuk pembelajaran sistem jangka panjang.
\end{enumerate}
\\ \hline
\end{tabular}
\end{table}

\begin{figure}[H]
    \centering
    % Ganti dengan nama file foto yang Anda simpan
%     \includegraphics[width=0.8\textwidth]{gambar/lampiran/validasi-prototipe-2025.jpg}
    \caption{Kegiatan validasi prototipe dan diskusi pengembangan lanjutan bersama Prof. Dr. apt. Maria Immaculata Iwo, M.Si.}
    \label{fig:validasi-prototipe}
\end{figure}